\begin{prob}{2.8}{}{関数論}
複素平面$\mathbb{C}$全体で正則な関数$f$が、任意の複素数$z$に対して
\[
|f(z)|\leqq 999\,|z|^{0.999}\,(\log(|z|+999))^{999}
\]
を満たすとき、$f(z)$は定数関数$0$であることを示せ。
\end{prob}

グルサの公式より, 任意の複素数$z$と正数$R>0$に対して
\[f'(z) = \int_{|\zeta - z| = R} \frac{f(\zeta)}{(\zeta - z)^2} \, d\zeta \]
これに絶対値を取って評価すれば
\begin{align*}
|f'(z)| &\leqq \int_{|\zeta - z| = R} \frac{|f(\zeta)|}{|\zeta - z|^2} \, |d\zeta| \\
&\leqq \int_{|\zeta - z| = R} \frac{999\,|\zeta|^{0.999}\,(\log(|\zeta|+999))^{999}}{|\zeta - z|^2} \, |d\zeta|
&\leqq \int_{|\zeta - z| = R} \frac{999\,(2R)^{0.999}\,(\log(2R+999))^{999}}{R^2} \, |d\zeta| \\
&= 999\,(2R)^{0.999}\,(\log(2R+999))^{999} \cdot \dfrac{2\pi R}{R^2} \\
&= 999\,(2R)^{0.999}\,(\log(2R+999))^{999} \cdot \dfrac{2\pi}{R} \\
&= 1998\pi \, R^{-0.001} (\log(2R+999))^{999}
&\to 0 \quad (R \to \infty)
\end{align*}
したがって$f'(z)=0$であるから, $f(z)$は定数関数である.
条件の不等式で$z=0$ｔとれば, $|f(0)|\leqq 0$であるから, $f(z)$は定数関数$0$である.
